\SourceLeaf{040}
\[
=fEx+gEy+hEz+\text{etc.}
\]
atque per \(Et\) denotabitur. Manifesto valor medius huius erroris fit \(=0\), sive error a parte constante liber erit. At valor medius quadrati \((Et)^2\), sive valor medius aggregati
\[
\begin{array}{llll}
ff(Ex)^2&+2fg\,Ex.Ey&+2fh\,Ex.Ez&+\text{etc.}\\
&+gg(Ey)^2&+2gh\,Ey.Ez&+\text{etc.}\\
&&+hh(Ez)^2&+\text{etc. etc.}
\end{array}
\]
per ea, quae in art. praec. exposuimus, aequalis fit producto ex \(mmp\) in aggregatum
\[
\begin{array}{llll}
ff[\alpha\alpha]&+2fg[\alpha\beta]&+2fh[\alpha\gamma]&+\text{etc.}\\
&+gg[\beta\beta]&+2gh[\beta\gamma]&+\text{etc.}\\
&&+hh[\gamma\gamma]&+\text{etc. etc.}
\end{array}
\]
sive producto ex \(mmp\) in valorem functionis \(\Omega-M\), qui prodit per substitutiones
\[
\xi=f,\quad\eta=g,\quad\zeta=h\quad\text{etc.}
\]
Denotando igitur hunc valorem determinatum functionis \(\Omega-M\) per \(\omega\), error medius metuendus, dum determinationi \(t=K\) adhaeremus, erit \(=m\BareRoot{p\omega}\), sive pondus huius determinationis \(=\frac1\omega\).

Quum indefinite habeatur
\[
\Omega-M=(x-A)\xi+(y-B)\eta+(z-C)\zeta+\text{etc.}
\]
patet, \(\omega\) quoque aequalem esse valori determinato expressionis
\[
(x-A)f+(y-B)g+(z-C)h+\text{etc.}
\]
sive valori determinato ipsius \(t-K\), qui prodit, si indeterminatis \(x,y,z\) etc. tribuuntur valores ii, qui respondent valoribus ipsarum \(\xi,\eta,\zeta\) etc. his \(f,g,h\) etc.

Denique observamus, si \(t\) indefinite in formam functionis ipsarum \(\xi,\eta,\zeta\) etc. redigatur, ipsius partem constantem necessario fieri \(=K\). Quodsi igitur indefinite fit
\[
t=F\xi+G\eta+H\zeta+\text{etc.}+K,\quad\text{erit}\quad\omega=fF+gG+hH+\text{etc.}
\]
\SourceLeaf{041}
\Article{30}
Functio \(\Omega\) valorem suum \textit{absolute minimum} \(M\), ut supra vidimus, nanciscitur, faciendo \(x=A,\allowbreak y=B,\allowbreak z=C\) etc., sive \(\xi=0,\allowbreak\eta=0,\allowbreak\zeta=0\) etc. Si vero alicui illarum quantitatum valor \textit{alius} iam tributus est, e. g. \(x=A+\Delta\), variantibus reliquis \(\Omega\) assequi potest valorem \textit{relative minimum}, qui manifesto obtinetur adiumento aequationum
\[
x=A+\Delta,\quad\frac{d\Omega}{dy}=0,\quad\frac{d\Omega}{dz}=0\quad\text{etc.}
\]
Fieri debet itaque \(\eta=0,\zeta=0\) etc., adeoque, quoniam
\[
x=A+[\alpha\alpha]\xi+[\alpha\beta]\eta+[\alpha\gamma]\zeta+\text{etc.},\quad\xi=\frac{\Delta}{[\alpha\alpha]}
\]
Simul habebitur
\[
y=B+\frac{[\alpha\beta]\Delta}{[\alpha\alpha]},\quad z=C+\frac{[\alpha\gamma]\Delta}{[\alpha\alpha]}\quad\text{etc.}
\]
Valor relative minimus ipsius \(\Omega\) autem fit \(=[\alpha\alpha]\xi\xi+M=M+\frac{\Delta\Delta}{[\alpha\alpha]}\). Vice versa hinc colligimus, si valor ipsius \(\Omega\) limitem praescriptum \(M+\mu\mu\) non superare debet, valorem ipsius \(x\) necessario inter limites \(A-\mu\BareRoot{[\alpha\alpha]}\) et \(A+\mu\BareRoot{[\alpha\alpha]}\) contentum esse debere. Notari meretur, \(\mu\BareRoot{[\alpha\alpha]}\) aequalem fieri errori medio in valore maxime plausibili ipsius \(x\) metuendo, si statuatur \(\mu=m\BareRoot{p}\), i. e. si \(\mu\) aequalis sit errori medio observationum talium, quibus pondus \(=1\) tribuitur.

Generalius investigemus valorem minimum ipsius \(\Omega\), qui pro valore dato ipsius \(t\) locum habere potest, denotante \(t\) ut in art. praec. functionem linearem \(fx+gy+hz+\text{etc.}+k\), et cuius valor maxime plausibilis \(=K\): valor praescriptus ipsius \(t\) denotetur per \(K+\kappa\). E theoria maximorum et minimorum constat, problematis solutionem petendam esse ex aequationibus
\[
\begin{aligned}
\frac{d\Omega}{dx}&=\theta\frac{dt}{dx}\\[3pt]
\frac{d\Omega}{dy}&=\theta\frac{dt}{dy}\\[3pt]
\frac{d\Omega}{dz}&=\theta\frac{dt}{dz}\quad\text{etc.}
\end{aligned}
\]
sive \(\xi=\theta f,\allowbreak\eta=\theta g,\allowbreak\zeta=\theta h\) etc., designante \(\theta\) multiplicatorem adhuc indeterminatum. Quare si, ut in art. praec., statuimus, esse \textit{indefinite}
\[
t=F\xi+G\eta+H\zeta+\text{etc.}+K
\]
habebimus
